\item Considere una interfaz horizontal entre el aire por encima de un vidrio de índice $1.55$.
\begin{enumerate}
    \item Trace un rayo de luz que incide desde el aire a un ángulo de incidencia de $30.0^\circ$. Determine los ángulos de los rayos reflejado y refractado e indíquelos en un diagrama.
    \item \textbf{¿Qué pasaría si?} Ahora suponga que el rayo de luz incide desde el vidrio a un ángulo de incidencia de $30.0^\circ$. Determine los ángulos de los rayos reflejado y refractado y señale los tres rayos en un nuevo diagrama.
    \item Para rayos que inciden desde el aire sobre la superficie de aire-vidrio, determine y tabule los ángulos de reflexión y refracción para todos los ángulos de incidencia a intervalos de $10.0^\circ$ desde $0^\circ$ hasta $90.0^\circ$.
    \item Haga lo mismo para rayos de luz que suben a la interfaz a través del vidrio.
\end{enumerate}